\documentclass[twoside]{article}

\usepackage[preprint]{aistats2027}
\usepackage[round,authoryear]{natbib}
\usepackage{amsmath,amssymb,amsthm}
\usepackage{mathtools}
\usepackage{microtype}
\usepackage{url}

\renewcommand{\bibname}{References}
\renewcommand{\bibsection}{\subsubsection*{\bibname}}
\bibliographystyle{apalike}

\newtheorem{theorem}{Theorem}
\newtheorem{proposition}[theorem]{Proposition}
\newtheorem{corollary}[theorem]{Corollary}
\newtheorem{lemma}[theorem]{Lemma}
\newcommand{\E}{\mathbb E}
\newcommand{\PP}{\mathbb P}
\newcommand{\cF}{\mathcal F}
\newcommand{\cX}{\mathcal X}
\newcommand{\TV}{\operatorname{TV}}
\newcommand{\err}{\operatorname{err}}
\newcommand{\Ber}{\operatorname{Ber}}
\newcommand{\Poi}{\operatorname{Poi}}
\newcommand{\1}{\mathbf 1}

\begin{document}
\runningauthor{}

\onecolumn
\aistatstitle{Supplement: When Sharing Fails Under Massart Noise}

\aistatsauthor{Pahan Dewasurendra}

\aistatsaddress{Johns Hopkins University}

\appendix

\section{External-Atom Moment Pairs}
\label{app:moment}

We give complete details for the moment construction. Let
\[
 d\sigma(u)=\frac{du}{\pi\sqrt{1-u^2}},\qquad u\in(-1,1),
\]
and let $T_j$ be the degree-$j$ Chebyshev polynomial of the first kind. The
orthonormal polynomials in $L_2(\sigma)$ are
\[
 \phi_0(u)=1,\qquad \phi_j(u)=\sqrt2T_j(u),\quad j\ge1.
\]
For $m\ge1$ and $z>1$, define
\[
 K_m(z)=\sum_{j=0}^m\phi_j(z)^2
 =1+2\sum_{j=1}^mT_j(z)^2.
\]

\begin{lemma}[External-atom representation]
\label{lem:app-moment}
There is a probability measure $\nu_m$ supported on $[-1,1]$ such that
\[
 \sigma_1=K_m(z)^{-1}\delta_z+
 (1-K_m(z)^{-1})\nu_m
\]
has the same moments as $\sigma$ through order $2m$. The atom mass
$K_m(z)^{-1}$ is maximal.
\end{lemma}

\begin{proof}
Let $w=K_m(z)^{-1}$. For a polynomial $P$ of degree at most $2m$, define
\[
 L(P)=\int_{-1}^1P(u)d\sigma(u)-wP(z).
\]
We first show that $L(P)\ge0$ whenever $P$ is nonnegative on $[-1,1]$.
The univariate Markov--Lukacs representation gives
\[
 P(u)=s(u)^2+(1-u^2)t(u)^2
\]
for polynomials satisfying $\deg(s)\le m$ and $\deg(t)\le m-1$. This
representation covers degree below $2m$ by allowing zero leading
coefficients. See \citet[Proposition 2.5.1]{wu2020polynomial}.

Expand $s=\sum_{j=0}^mc_j\phi_j$. Cauchy--Schwarz gives
\begin{align}
 s(z)^2
 &=\left(\sum_{j=0}^mc_j\phi_j(z)\right)^2\notag\\
 &\le\left(\sum_{j=0}^mc_j^2\right)
 \left(\sum_{j=0}^m\phi_j(z)^2\right)
 =K_m(z)\int s^2d\sigma.
 \label{eq:app-kernel-cs}
\end{align}
Since $wK_m(z)=1$ and $z>1$,
\begin{align*}
 L(P)
 &=\int s^2d\sigma-ws(z)^2
 +\int(1-u^2)t(u)^2d\sigma\\
 &\quad+w(z^2-1)t(z)^2\ge0.
\end{align*}
The truncated Riesz--Haviland theorem
\citep[Theorem 2.5.2]{wu2020polynomial} now gives a nonnegative measure
$\tau$ supported on $[-1,1]$ such that
\[
 L(P)=\int P\,d\tau
\]
for every polynomial of degree at most $2m$. Taking $P=1$ shows that
$\tau([-1,1])=1-w$. Since $K_m(z)>1$, normalizing
$\nu_m=\tau/(1-w)$ proves existence and moment equality.

For maximality, suppose a probability measure with the same first $2m$
moments has an atom of mass $w'$ at $z$ and all remaining mass on
$[-1,1]$. Every degree-$m$ polynomial $s$ then satisfies
\[
 w's(z)^2\le\int s^2d\sigma.
\]
Take $s_z(u)=\sum_{j=0}^m\phi_j(z)\phi_j(u)$. Orthonormality gives
$s_z(z)=K_m(z)$ and $\int s_z^2d\sigma=K_m(z)$. Therefore
$w'\le1/K_m(z)$.
\end{proof}

\begin{lemma}[Closed form]
\label{lem:kernel-form}
If $A=\operatorname{arcosh}(z)$, then
\[
 K_m(z)=m+\frac12+\frac{\sinh((2m+1)A)}{2\sinh A}.
\]
For every fixed $z>1$, there are positive constants $c_z,C_z$ such that
\[
 c_ze^{2mA}\le K_m(z)\le C_ze^{2mA}
\]
for all $m\ge1$.
\end{lemma}

\begin{proof}
Since $T_j(z)=\cosh(jA)$,
\begin{align*}
 K_m(z)
 &=1+2\sum_{j=1}^m\cosh^2(jA)\\
 &=m+1+\sum_{j=1}^m\cosh(2jA)\\
 &=m+\frac12+\frac{\sinh((2m+1)A)}{2\sinh A}.
\end{align*}
The exponential comparison follows because $A>0$ is fixed.
\end{proof}

For the learning construction, apply the affine map
\[
 q=\frac{\eta}{2}(u+1).
\]
It sends $[-1,1]$ to $[0,\eta]$. The point $q=1-\eta$ corresponds to
\[
 z_\eta=\frac{2(1-\eta)}{\eta}-1>1.
\]
Let $\mu_0$ be the image of $\sigma$. Lemma~\ref{lem:app-moment} gives a
probability measure $\widetilde\nu_m$ supported on $[0,\eta]$ such that
\begin{equation}
 \mu_1=w_m\delta_{1-\eta}+(1-w_m)\widetilde\nu_m,
 \qquad w_m=K_m(z_\eta)^{-1},
 \label{eq:app-bias-pair}
\end{equation}
matches moments through order $2m$ with $\mu_0$.

\section{The Single-Source Experiment}
\label{app:testing}

Fix $d\ge1$ and $\alpha\in(0,1]$. The covariate domain is
$\{0,1,\ldots,d\}$ with
\[
 p_X(0)=1-\alpha,
 \qquad p_X(x)=\alpha/d\quad(x\in[d]).
\]
Set the bias at coordinate zero to any common value $q_0\in[0,\eta]$.
Under hypothesis $h\in\{0,1\}$, first draw
$q_1,\ldots,q_d$ independently from $\mu_h$. Keep this vector fixed, then
draw each label from $Y\mid X=x\sim\Ber(q_x)$.

\subsection{Poissonization}

Let the number of observations be Poisson with mean $N$. For one light
coordinate, write $S$ and $T$ for its zero-label and one-label counts. Given
$q$, Poisson splitting gives
\[
 S\sim\Poi(\lambda(1-q)),
 \qquad T\sim\Poi(\lambda q),
 \qquad \lambda=N\alpha/d,
\]
independently. The mixed probability is
\begin{equation}
 Q_h(s,t)=\frac{e^{-\lambda}\lambda^{s+t}}{s!t!}
 \E_{q\sim\mu_h}[q^t(1-q)^s].
 \label{eq:app-count-law}
\end{equation}

\begin{lemma}[Poissonized total variation]
\label{lem:app-tv}
For the complete Poissonized observations,
\[
 \TV(\mathsf P_0^N,\mathsf P_1^N)
 \le d\,\PP(\Poi(\lambda)\ge2m+1).
\]
In particular, if
\begin{equation}
 N\le N_0:=\frac{2m+1}{4e\alpha}
 d^{\frac{2m}{2m+1}},
 \label{eq:app-N0}
\end{equation}
then this total variation is at most $1/64$.
\end{lemma}

\begin{proof}
For $s+t\le2m$, the expectation in \eqref{eq:app-count-law} is the
expectation of a polynomial of degree at most $2m$. It is equal under the two
moment-matching priors. Under either prior, $S+T\sim\Poi(\lambda)$ because its
law does not depend on $q$. Hence the one-coordinate distance is at most
$\PP(\Poi(\lambda)\ge2m+1)$.

After integrating out the independent latent biases, the light-coordinate
count pairs are independent. Product subadditivity gives the first claim. If
$N\le N_0$, then
\[
 \lambda\le\frac{2m+1}{4e}d^{-1/(2m+1)}.
\]
The Poisson Chernoff bound yields
\begin{align*}
 d\,\PP(\Poi(\lambda)\ge2m+1)
 &\le d\left(\frac{e\lambda}{2m+1}\right)^{2m+1}\\
 &\le4^{-(2m+1)}\le\frac1{64}.
\end{align*}
The heavy-coordinate process is identical under both priors. Conditional on
all counts, the ordering of observations is generated by a common Markov
kernel. Thus the same bound applies to the full ordered data.
\end{proof}

\subsection{A fixed-budget lower bound}

\begin{lemma}[De-Poissonization]
\label{lem:depoisson}
Assume $N_0\ge48$. Every test whose two error probabilities are at most $1/5$
uses more than $N_0/4$ fixed-budget observations.
\end{lemma}

\begin{proof}
Suppose a test uses at most $n\le N_0/4$ observations. In a Poisson experiment
with mean $2\max\{n,12\}$, run the test on the first $n$ observations when
they are available. Otherwise output a fair coin. The Poisson mean is at most
$N_0$. A multiplicative lower-tail bound gives
\[
 \PP(\Poi(2\max\{n,12\})\ge n)\ge0.9.
\]
The Poissonized test has success probability at least
\[
 \frac12+0.9\left(\frac45-\frac12\right)>\frac34
\]
under both priors. Two distributions that admit such a test have total
variation above $1/2$. This contradicts Lemma~\ref{lem:app-tv}.
\end{proof}

Combining \eqref{eq:app-N0} and Lemma~\ref{lem:depoisson}, every test requires
at least
\begin{equation}
 L_m:=\frac{2m+1}{16e\alpha}
 d^{\frac{2m}{2m+1}}
 \label{eq:app-Lm}
\end{equation}
observations, once this quantity is larger than a universal constant.

\section{Valid Massart Alternatives}
\label{app:gap}

Let $\gamma=1-2\eta$ and set
\begin{equation}
 \alpha=\frac{16\epsilon}{\gamma w_m}.
 \label{eq:app-alpha}
\end{equation}
The first condition of the main theorem, with a sufficiently small universal
constant, ensures $\alpha\le1$.

Under the null, every light bias lies in $[0,\eta]$, so the zero function is
Bayes optimal. Under the alternative, define
\[
 H=\sum_{x=1}^d\1\{q_x=1-\eta\}.
\]
Then $H\sim\operatorname{Bin}(d,w_m)$. Let
\begin{equation}
 \mathcal A=\{dw_m/2\le H\le2dw_m\}.
 \label{eq:app-A}
\end{equation}

\begin{lemma}[Gap event]
\label{lem:gap-event}
If $dw_m\ge64$, then $\PP(\mathcal A)\ge0.99$. On $\mathcal A$, the Bayes
classifier is the indicator of the high coordinates, and the excess error of
zero satisfies
\[
 8\epsilon\le\Delta_0\le32\epsilon.
\]
Its disagreement probability with Bayes is at most $32\epsilon/\gamma$.
\end{lemma}

\begin{proof}
The two-sided multiplicative Chernoff bound gives
\[
 \PP(\mathcal A^c)
 \le e^{-dw_m/8}+e^{-dw_m/3}<0.01.
\]
The only points where zero differs from Bayes are the $H$ high coordinates.
Every such point has label margin $2(1-\eta)-1=\gamma$. Therefore
\[
 \Delta_0=\frac{\alpha\gamma H}{d}.
\]
Substituting \eqref{eq:app-alpha} and the bounds in
\eqref{eq:app-A} proves the excess interval. The disagreement mass is
$\alpha H/d\le2\alpha w_m=32\epsilon/\gamma$.
\end{proof}

Substitution into \eqref{eq:app-Lm} gives
\begin{equation}
 L_m=\frac{(2m+1)\gamma w_m}{256e\epsilon}
 d^{\frac{2m}{2m+1}}.
 \label{eq:app-Lm-final}
\end{equation}

\section{The Variance-Sensitive Holdout}
\label{app:holdout}

The next result is stated conditionally on any fixed valid environment and
any fixed output predictor. It therefore remains valid after conditioning on
the learner's transcript.

\begin{lemma}[Holdout test]
\label{lem:holdout}
Let $f^\star$ be a Bayes classifier satisfying the Massart condition with
margin lower bound $\gamma$. Let $\widehat f$ have excess error at most
$\epsilon$, and let $f_0$ be zero. Under a null environment where
$f_0=f^\star$, define
\[
 D=\err(f_0)-\err(\widehat f).
\]
Then $D\in[-\epsilon,0]$. Under an alternative in $\mathcal A$,
$D\ge7\epsilon$. In both cases, the empirical difference from
\begin{equation}
 M=\left\lceil
 \frac{20\log(2/\delta_v)}{\gamma\epsilon}
 \right\rceil
 \label{eq:app-M}
\end{equation}
fresh samples is within $2\epsilon$ of $D$ with probability at least
$1-\delta_v$. Hence thresholding it at $3\epsilon$ distinguishes the cases.
\end{lemma}

\begin{proof}
For every predictor $f$, the pointwise Massart identity gives
\[
 \err(f)-\err(f^\star)
 =\E[(1-2\eta^\star(X))\1\{f(X)\ne f^\star(X)\}],
\]
where $\eta^\star(X)\le\eta$. Thus an $\epsilon$-optimal predictor obeys
\begin{equation}
 \PP(\widehat f(X)\ne f^\star(X))\le\epsilon/\gamma.
 \label{eq:app-disagree}
\end{equation}
Under the null, $f_0=f^\star$, which proves $D\in[-\epsilon,0]$. Under an
alternative in $\mathcal A$,
\[
 D=\Delta_0-(\err(\widehat f)-\err(f^\star))\ge7\epsilon.
\]

For a fresh sample let
\[
 Z=\1\{f_0(X)\ne Y\}-\1\{\widehat f(X)\ne Y\}.
\]
It vanishes whenever $f_0(X)=\widehat f(X)$. Under the null,
$\operatorname{Var}(Z)\le\epsilon/\gamma$. Under the alternative, Lemma
\ref{lem:gap-event} and \eqref{eq:app-disagree} give
\[
 \operatorname{Var}(Z)
 \le\PP(f_0(X)\ne\widehat f(X))
 \le33\epsilon/\gamma.
\]
Also $|Z-\E Z|\le2$. Bernstein's inequality yields
\[
 \PP(|\overline Z-D|>2\epsilon)
 \le2\exp(-M\gamma\epsilon/20)\le\delta_v.
\]
The two claimed intervals are separated by the threshold $3\epsilon$.
\end{proof}

\section{Adaptive Multi-Source Lifting}
\label{app:lifting}

Let the covariate space be the disjoint union of $k$ blocks
$\cX_1,\ldots,\cX_k$. Block $i$ has one heavy point $x_{i,0}$ and light
points $x_{i,1},\ldots,x_{i,d}$. Define
\begin{equation}
 \cF=\{f_0\}\cup
 \bigcup_{i=1}^k\{\1_S:S\subseteq\{x_{i,1},\ldots,x_{i,d}\}\}.
 \label{eq:app-class}
\end{equation}

\begin{lemma}[VC dimension]
\label{lem:vc}
The class in \eqref{eq:app-class} has VC dimension exactly $d$.
\end{lemma}

\begin{proof}
The $d$ light points in any one block are shattered. A set containing a heavy
point cannot be shattered because every function labels that point zero. A
set of $d+1$ light points must span at least two blocks. The labeling that is
one on a point from each of two blocks is not realized by $\cF$.
\end{proof}

Source $i$ is supported on $\cX_i$ with mass $1-\alpha$ on $x_{i,0}$ and
$\alpha/d$ on each light point. Set every heavy bias to a fixed value in
$[0,\eta]$.

Under the common null prior $\Pi_0$, independently draw every source's light
bias vector from $\mu_0^d$. Each realization is a valid Massart instance with
shared Bayes classifier $f_0$. For a fixed $i$, prior $\Pi_i$ replaces only
source $i$ by $\mu_1^d$. Its high-coordinate indicator is in $\cF$ and equals
zero on all other supports. It is therefore a shared Bayes classifier for
every realization under $\Pi_i$.

\begin{theorem}[Adaptive lifting]
\label{thm:app-lifting}
Suppose a personalized learner succeeds with probability at least $0.9$ on
every valid instance. If
\begin{equation}
 m\gamma^2w_m d^{\frac{2m}{2m+1}}\ge C
 \label{eq:app-domination}
\end{equation}
for a sufficiently large universal $C$, then under the common null
\[
 \E_{\Pi_0}[T_i]\ge L_m/100
\]
for every source $i$, where $T_i$ is its query count and $L_m$ is given by
\eqref{eq:app-Lm-final}.
\end{theorem}

\begin{proof}
Fix $i$ and suppose instead that $\E_{\Pi_0}[T_i]<L_m/100$. We build a test
between the single-source priors from Appendix~\ref{app:testing}. Run the MDL
learner. Whenever it queries source $i$, return a sample from the unknown
testing oracle after shifting its covariate into block $i$. For each other
source, independently draw a null bias vector once and simulate samples from
that fixed environment.

If the source-$i$ query count reaches $L_m/2$, stop and output alternative. If
the learner terminates first, take the source-$i$ output and run the holdout
test from Lemma~\ref{lem:holdout} with $\delta_v=0.02$.

Under the null, Markov's inequality gives
\[
 \PP_{\Pi_0}(T_i\ge L_m/2)<0.02.
\]
The learner succeeds with probability at least $0.9$, including after
averaging over the valid randomized null. Conditional on its transcript and
success, the holdout accepts with probability at least $0.98$. A union bound
shows that the constructed test accepts the null with probability above
$0.8$.

Under the alternative, truncation is a correct output. If truncation does not
occur, the intersection of $\mathcal A$, learner success, and holdout success
makes the test output alternative. This intersection has probability at
least $0.99-0.1-0.02>0.8$. Pointwise correctness of the learner applies to
every alternative realization, not only after conditioning on $\mathcal A$.

The unknown-oracle budget is at most $L_m/2+M$. From
\eqref{eq:app-Lm-final} and \eqref{eq:app-M}, condition
\eqref{eq:app-domination} with a sufficiently large $C$ ensures
$M\le L_m/4$. It also ensures $N_0=4L_m\ge48$. We have constructed a test
using less than $L_m$ observations with both errors below $1/5$, contradicting
Lemma~\ref{lem:depoisson}. Therefore $\E_{\Pi_0}[T_i]\ge L_m/100$.
\end{proof}

Summing Theorem~\ref{thm:app-lifting} over $i$ under the same null gives
\begin{align*}
 \E_{\Pi_0}\left[\sum_{i=1}^kT_i\right]
 &\ge\frac{kL_m}{100}\\
 &=\Omega\!\left(
 \frac{k m\gamma w_m}{\epsilon}
 d^{\frac{2m}{2m+1}}
 \right).
\end{align*}
Since every realization in the support of $\Pi_0$ is valid, some fixed
environment has at least this expected cost. This proves the main theorem for
expected worst-case sample complexity. A deterministic worst-case sample cap
is at least the same quantity.

\section{Optimization and Comparison}
\label{app:optimization}

\subsection{Fixed exponent loss}

Fix $\eta\in(0,1/2)$ and $\zeta>0$. Choose a fixed integer $m$ satisfying
\[
 \frac1{2m+1}\le\zeta.
\]
Then $w_m>0$ is a constant depending only on $\eta$ and $\zeta$. Choose
$\epsilon\le c\gamma w_m$. For all sufficiently large $d$, the count and
holdout conditions hold. The main theorem gives
\[
 \Omega_{\eta,\zeta}\!\left(
 \frac{k}{\epsilon}d^{2m/(2m+1)}
 \right)
 =\Omega_{\eta,\zeta}\!\left(
 \frac{k}{\epsilon}d^{1-\zeta}
 \right).
\]

\subsection{Polynomial accuracy}

Fix $\beta>0$ and suppose $\epsilon\le d^{-\beta}$. Write
$D=\log d$ and $A=A_\eta>0$. Choose an integer $m$ such that
\begin{equation}
 2m\asymp\sqrt{D/A}.
 \label{eq:app-m-choice}
\end{equation}
For all sufficiently large $d$, Lemma~\ref{lem:kernel-form} gives
\[
 w_m\ge c_\eta\exp(-2mA).
\]
The accuracy condition holds because
\[
 \log(1/\epsilon)\ge\beta D
 \gg2mA+O_\eta(1)=O_\eta(\sqrt D).
\]
Similarly,
\[
 dw_m\ge c_\eta\exp(D-O_\eta(\sqrt D))\longrightarrow\infty.
\]
For the holdout condition,
\begin{align*}
 m\gamma^2w_m d^{2m/(2m+1)}
 &\ge c_\eta m\exp\left(
 D-2mA-\frac{D}{2m+1}
 \right)\\
 &\ge\exp(D-O_\eta(\sqrt D))\longrightarrow\infty.
\end{align*}
Thus all conditions hold. The main lower bound is at least
\begin{align*}
 \frac{c_\eta kdm}{\epsilon}
 \exp\left(-2mA-\frac{D}{2m+1}\right)
 &\ge\frac{c_{\eta,\beta}kd}{\epsilon}
 \exp(-C_\eta\sqrt{\log d}).
\end{align*}

\subsection{Separate empirical risk minimization}

For completeness, the Massart identity used in
\eqref{eq:app-disagree} also gives the Bernstein condition
\[
 \operatorname{Var}(\ell_f-\ell_{f^\star})
 \le\frac{\err(f)-\err(f^\star)}{\gamma}.
\]
Standard localized VC empirical risk minimization therefore learns one
source to excess $\epsilon$ and failure probability $\delta/k$ from
\[
 O\!\left(
 \frac{d\log(1/\epsilon)+\log(k/\delta)}{\gamma\epsilon}
 \right)
\]
samples \citep{boucheron2005classification,massart2006risk}. Applying it
independently to all sources and taking a union bound gives the upper
benchmark stated in the main paper.

\bibliography{references}

\end{document}
